\section{Query-efficient quantum algorithm for minimum cut} \label{sec:query-efficient}
We first describe a query-efficient quantum algorithm to find the weight and shores of a minimum cut.
In \cref{sec:time-efficient} we make this algorithm time-efficient.
Our quantum query algorithm for minimum cut mainly relies on \cref{thm:RSWweighted}, and is inspired by a classical randomized algorithm for edge connectivity in the cut query model by Rubinstein, Schramm, and Weinberg (RSW) \cite{RSW18}.  
The RSW cut query algorithm is based on 4 subroutines whose input/output behavior we describe in Algorithms 1--4 below.  
For weighted graphs, we need an additional subroutine to compute the maximum weight of an edge in the graph which is stated in \cref{alg:maxweight}.
We describe all these subroutines in an abstract way to make it easy to (i) describe the time-efficient algorithm in the next section, and (ii) to instantiate this algorithm for other query models in the future.
We indicate oracle access to $G$ by square brackets and put the parameters explicitly given to the routines in parentheses.

\begin{algorithm}[H]
\caption{FindMinStar$[G](\delta)$}
\label{alg:mindeg}
 \hspace*{\algorithmicindent} \textbf{Input:} Oracle access to a weighted graph $G$, error parameter $\delta$. \\
 \hspace*{\algorithmicindent} \textbf{Output:} With probability at least $1-\delta$ output $v \in \argmin_{u \in V} w(\Delta_G(\{u\}))$ and 
 $d_{\min} = \min_{u \in V} w(\Delta_G(\{u\}))$.
\end{algorithm}

\begin{algorithm}[H]
\caption{Cut-Sparsifier$[G](\varepsilon, \delta)$}
\label{alg:sparsify}
 \hspace*{\algorithmicindent} \textbf{Input:} Oracle access to a weighted graph $G$, sparsifier accuracy parameter $\varepsilon$, error parameter $\delta$. \\
 \hspace*{\algorithmicindent} \textbf{Output:} With probability at least $1-\delta$ output an integer-weighted $\varepsilon$-cut sparsifier $H$ of $G$ with $\tO(n/\epsilon^2)$ edges.
\end{algorithm}

\begin{algorithm}[H]
\caption{LearnCutAtoms$(H,\lambda, \delta)$}
 \hspace*{\algorithmicindent} \textbf{Input:} Adjacency array description of $H$, cut threshold $\lambda$, and error parameter $\delta$. \\
 \hspace*{\algorithmicindent} \textbf{Output:}
 Define the set $\Tcal = \{ X: |X|, |\overline{X}| \ge 2, w(\Delta_H(X)) \le \lambda \}$.
 With probability at least $1-\delta$ output $\atoms(\Tcal)$.
\end{algorithm}

\begin{algorithm}[H]
\caption{LearnContraction$[G](\Pcal, M, \delta)$}
\label{alg:contract}
 \hspace*{\algorithmicindent} \textbf{Input:} Oracle access to a weighted graph $G$, a partition $\Pcal$ of $V(G)$, a natural number $M$, and error parameter $\delta$. \\
 \hspace*{\algorithmicindent} \textbf{Output:} Let $G' = \Contract(G,\Pcal)$.  With probability at least $1-\delta$ return $G'$ if the number of edges of $G'$ is at most $M$, 
 and otherwise return $\mathrm{NULL}$.
\end{algorithm}

\begin{algorithm}[H]
\caption{FindMaxWeight$[G](\delta)$}
\label{alg:maxweight}
 \hspace*{\algorithmicindent} \textbf{Input:} Oracle access to a weighted graph $G$, error parameter $\delta$. \\
 \hspace*{\algorithmicindent} \textbf{Output:} With probability at least $1-\delta$ output $\tau$, the maximum weight of an edge of $G$.
\end{algorithm}

We combine these subroutines in \cref{alg:mincut} to give a template for solving the minimum cut problem in an abstract query model.
\begin{algorithm}[H]
\caption{Query algorithm for minimum cut}
\label{alg:mincut}
 \hspace*{\algorithmicindent} \textbf{Input:} Oracle access to a weighted graph $G$ \\
 \hspace*{\algorithmicindent} \textbf{Output:} $\lambda(G)$ and the shores of a minimum cut of $G$.
\begin{algorithmic}[1]
\State $(v,d_{\min}) \gets$ FindMinStar$[G](\frac{1}{20})$.
\State $\tau \gets$ FindMaxWeight$[G](\frac{1}{20})$.
\State $H = (V,w') \gets$ Cut-Sparsifier$[G](\frac{1}{100}, \frac{1}{20})$.
\State Compute $\lambda(H)$.
\State $\Pcal = \{ S_1,\dots,S_k \} \gets$ LearnCutAtoms$(H,(1+\frac{1}{100})\lambda(H),\frac{1}{20})$.
\State $G' \gets$ LearnContraction$[G](\Pcal, 100 \tau n, \frac{1}{20})$.  If $G' = \mathrm{NULL}$ then abort.
\State Compute the weight $\lambda(G')$ and shores $(Y,V(G') \setminus Y)$ of a minimum cut in $G'$.
\State If $d_{\min} \le \lambda(G')$ output $(d_{\min},(\{v\},V\backslash\{v\}))$. Otherwise, let $Z = \cup_{S_i \in Y} S_i$ and output $(\lambda(G'), (Z, \overline{Z}))$.
\end{algorithmic}
\end{algorithm}

\begin{theorem}
\label{thm:rsw_alg}
Let $G$ be a weighted graph with $n$ vertices, minimum edge weight at least 1, and maximum edge weight $\tau$. 
\cref{alg:mincut} finds the weight and shores of a minimum cut of $G$ with probability at least $3/4$.  The number of queries of the algorithm is the sum of the number of queries 
of the subroutines FindMinStar$[G](\frac{1}{20})$, FindMaxWeight$[G](\frac{1}{20})$, Cut-Sparsifier$[G](\frac{1}{100}, \frac{1}{20})$, and LearnContraction$[G](\Pcal, 100 \tau n, \frac{1}{20})$.
\end{theorem}

\begin{proof}
Queries to the input graph $G$ are only made in steps $1,2,3,$ and $6$. This gives the statement about the complexity of the algorithm.  

Next let us deal with the error probability.  With probability at least $16/20$ steps 1--5 return correctly by the definition of these subroutines and 
the error parameter provided.  Let us now assume this is the case.   Then $H =(V,w')$ is a valid $\varepsilon$-sparsifier of $G$ for $\varepsilon = 1/100$.  
Let $X \in \Tcal$.  Then we have 
\[
w(\Delta_G(X)) \le (1+\varepsilon) w'(\Delta_H(X)) \le (1+\varepsilon)(1+3\varepsilon)\lambda(H) \le (1+\varepsilon)^2(1+3\varepsilon)\lambda(G) \enspace .
\]
We have $(1+\varepsilon)^2(1+3\varepsilon) \le \frac{11}{10}$ by the choice of $\varepsilon$, and so $w(\Delta_G(X)) \leq \frac{11}{10} \lambda(G) \leq \lambda(G) + \frac{1}{10} d_{\min}$ 
since $\lambda(G) \le d_{\min}$.  Thus by \cref{thm:RSWweighted}, the total weight of edges in $\Contract(G, \Pcal)$ will be at most $100\tau n$.
As we assume the minimum weight of an edge is at least $1$, the number of edges in $\Contract(G, \Pcal)$ will also be at most $100\tau n$. 
Hence except with probability at most $1/20$, LearnContraction will correctly return $\Contract(G, \Pcal)$ in step~5.

We have now argued that with probability at least $3/4$ all subroutines will correctly return.  We now argue correctness assuming that this is the case.  
In this case, $G'$ will be a valid contraction of $G$ and so $\lambda(G') \ge \lambda(G)$.  Thus if $\lambda(G)$ is achieved by a star cut the 
algorithm will return correctly.  

Let us now assume that $d_{\min} > \lambda(G)$ and let $\Delta_G(X)$ be a non-star cut with $w(\Delta_G(X)) = \lambda(G)$.  We have 
\[
w'(\Delta_H(X)) \le (1+\varepsilon) w(\Delta_G(X)) = (1+\varepsilon) \lambda(G) \le \frac{1+\varepsilon}{1-\varepsilon} \lambda(H) \le (1+3\varepsilon) \lambda(H)\enspace ,
\]
where the last step holds as $\varepsilon \le \frac{1}{3}$.  This means $X \in \Tcal$ and therefore no edge of $\Delta_G(X)$ will be 
contracted in $G' = \Contract(G,\Pcal)$.  Thus $\lambda(G') \le \lambda(G)$ and as the edge connectivity cannot decrease in a contraction in fact $\lambda(G') = \lambda(G)$.
Hence the algorithm returns correctly in step~8.
\end{proof}

\begin{lemma} \label{lem:compl-subroutines}
Let $G = (V,w)$ be a weighted graph with $n$ vertices and $m$ edges.
Subroutines FindMinStar$[G](\frac{1}{20})$, FindMaxWeight$[G](\frac{1}{20})$, Cut-Sparsifier$[G](\frac{1}{100}, \frac{1}{20})$ can be implemented by a quantum algorithm 
with query and time complexity  $\tO(n^{3/2})$ in the adjacency matrix and $\tO(\sqrt{mn})$ in the adjacency array model.  

LearnContraction$[G](\Pcal, 100 \tau n, \frac{1}{20})$ can be implemented by a quantum algorithm with query and time complexity  
$\tO(n^{3/2}\sqrt{\tau})$ in the adjacency matrix and $\tO(\sqrt{mn\tau})$ in the adjacency array model.
\end{lemma}

\begin{proof}
First note that in the adjacency array model we may assume that $m \ge n$. Otherwise, $\sqrt{mn} \ge m$ and we can perform each task classically in $\tO(m)$ time and queries.
We consider each of the subroutines in turn:

\textbf{FindMinStar$[G](\frac{1}{20})$:} In the adjacency matrix model we can compute $w(\Delta_G(\{v\})$ with $n-1$ classical queries to the adjacency matrix.  
We can compose this with quantum minimum finding to find the minimum weight of a star cut and a vertex realizing this in query and time complexity $\tO(n^{3/2})$ by \cref{thm:qmin}.

In the adjacency array model we first classically query the degrees of all the vertices with $n$ queries.  In a simple graph this suffices to determine the minimum weight of a star cut.
In a weighted graph we continue as follows.  For $1 \leq \ell \leq \ceil{\log n}$, define the bucket $B_\ell \subseteq V$ as the subset of nodes $v$ that have degree in $[2^{\ell-1}, 2^\ell)$.
As the sum of the degrees is $2m$ we have that $|B_\ell| \leq 2m/2^{\ell-1}$.
Finding the minimum $\min_{v \in B_\ell} w(\Delta_G(\{v\}))$ over a single bucket has quantum query and time complexity $\tO(\sqrt{mn})$: we can compute $w(\Delta_G(\{v\}))$ for a single $v \in B_\ell$ using at most $2^\ell$ classical queries, and 
then do quantum minimum finding over the $|B_\ell| \leq 2m/2^{\ell-1}$ nodes in $B_\ell$.
This has total query and time complexity $\tO(2^\ell \sqrt{2m/2^{\ell-1}}) \in \tO(\sqrt{m 2^\ell}) \in \tO(\sqrt{mn})$.
We do this for each of the $\ceil{\log n}$ buckets and we output the minimum overall weight and a vertex realizing this.
This yields a total time and query complexity $\tO(\sqrt{mn})$.

\textbf{FindMaxWeight$[G](\frac{1}{20})$:}
This amounts to finding the maximum of a set of $n^2$ numbers in the adjacency matrix model, or $m$ numbers in the adjacency list model.
By \cref{thm:qmin} this has query and time complexity $\tO(n)$ and $\tO(\sqrt{m})$, respectively.

\textbf{Cut-Sparsifier$[G](\frac{1}{100}, \frac{1}{20})$:} A $\frac{1}{100}$-cut sparsifier with $\tO(n/\epsilon^2)$ edges can be constructed with high probability in query and time complexity $\tilde O(n^{3/2})$ in the adjacency matrix model or $\tilde O(\sqrt{mn})$ in the adjacency array model by \cref{thm:qsparsifier}.  

\textbf{LearnContraction$[G](\Pcal, 100\tau n, \frac{1}{20})$:} 
First we handle a trivial case.  If $\tau \ge n$ then we can classically learn the input in time $n^2 = O(n^{3/2}\sqrt{\tau})$ in the adjacency matrix model and 
time $m = O(\sqrt{mn\tau})$ in the adjacency array model.  Thus we can assume $\tau < n$.  

First we do the adjacency matrix case.  Let $x \in \R^{\binom{n}{2}}$ be a vector whose entries are labeled by elements of $V^{(2)}$ and where $x(e) = w(e)$ if the endpoints of $e$ are in distinct elements of $\Pcal$ and $x(e) = 0$ otherwise.  
A query to an entry of $x$ can be answered with one query to the adjacency matrix of $G$.  Let $\hat{x} \in \{0,1\}^{\binom{n}{2}}$ be defined by $\hat{x}(e) = 1$ if $w(e) > 0$ and $\hat{x}(e) = 0$ otherwise.  We can also answer a query to $\hat{x}$ 
with one query to the adjacency matrix of $G$.  By \cref{thm:qsearch} in query and time complexity 
$\tO(n^{3/2}\sqrt{\tau})$ in the adjacency matrix model we can with probability at least $9/10$ 
output $\hat{x}$ if $|\hat{x}| \le 100 \tau n$ and otherwise output NULL.  We can then classically query $x$ in the non-zero locations of $\hat{x}$ with $100\tau n =O(n^{3/2} \sqrt{\tau})$ more classical queries to output $x$.  
This fulfils the specification of LearnContraction.

Similarly, in the adjacency array model let $x \in \R^{m}$ be labeled by entries of the adjacency array of $G$ and define $x(e) = w(e)$ if the endpoints of $e$ are in distinct elements of $\Pcal$ and 
$x(e) = 0$ otherwise.  Let $\hat{x}(e)=1$ if $x(e) >0$ and $\hat{x}(e)=0$ otherwise as before.
A query to an entry of $x$ or $\hat{x}$ can be answered with one query to the adjacency array of $G$.  Again by \cref{thm:qsearch}, in query and time complexity $\tO(\sqrt{mn \tau})$ in the adjacency array model we can with probability at least $9/10$ 
output $\hat{x}$ if $|\hat{x}| \le 100 \tau n$ and otherwise output NULL.  With $100 \tau n = O(\sqrt{mn\tau})$ more queries we can then output $x$.
\end{proof}

